%?====== Configuración del documento ===============================
\documentclass[12pt,letterpaper]{article}
\usepackage[spanish]{babel}
\renewcommand{\familydefault}{\sfdefault} % Usa Arial/Helvetica
\usepackage{helvet}
\renewcommand{\normalsize}{\fontsize{11pt}{13pt}\selectfont}
\usepackage{geometry}
\geometry{margin=2.5cm}
\usepackage{setspace}
\onehalfspacing % Interlineado 1.15 approx / 1.5
\setlength{\parindent}{0pt}
\setlength{\parskip}{0.1em}
\usepackage{graphicx}
%?===================================================================

%?============= Encabezados ===================
\usepackage{fancyhdr}
\pagestyle{fancy}
\fancyhf{}
\lhead{Ejercicio 2}          % Izquierda
\chead{\today}               % Centro
\rhead{Flores Olivares Omar} % Derecha

% Línea divisora bajo el encabezado
\renewcommand{\headrulewidth}{0.4pt}
\setlength{\headheight}{14.5pt}
%?=============================================

\title{\textbf{Ejercicio 2. El Sistema Gestor en un Contenedor: Docker}}
\author{Omar Flores Olivares}
\date{\today}

\begin{document}
\maketitle

\section*{Parte A. Investigación}

\subsection*{1. Concepto de Contenedor y Diferencias con una Máquina Virtual}
Un \textbf{contenedor} es un entorno ejecutable liviano, aislado y portátil que empaqueta una aplicación junto con todas sus dependencias, librerías y archivos de configuración requeridos para su correcto funcionamiento. A diferencia de una solución de virtualización tradicional, la diferencia fundamental radica en la arquitectura subyacente:

\begin{itemize}
    \item \textbf{Arranque:} Las máquinas virtuales (VM) requieren iniciar un sistema operativo huésped completo (\textit{Guest OS}) sobre un hipervisor, lo que demora minutos. Los contenedores comparten el núcleo (\textit{kernel}) del sistema operativo anfitrión y se inician en cuestión de milisegundos a pocos segundos.
    \item \textbf{Tamaño:} Una máquina virtual incluye una imagen completa del sistema operativo con sus propios controladores y servicios, pesando habitualmente varios gigabytes (GB). Una imagen de contenedor contiene únicamente la aplicación y sus dependencias de usuario, ocupando únicamente decenas o cientos de megabytes (MB).
    \item \textbf{Aislamiento:} Una máquina virtual ofrece un aislamiento fuerte a nivel de hardware garantizado por el hipervisor. Los contenedores proporcionan aislamiento a nivel de proceso utilizando funcionalidades del \textit{kernel} de Linux (como \textit{namespaces} y \textit{cgroups}), lo que reduce el sobrecosto de recursos procesales.
\end{itemize}

\subsection*{2. Definiciones}
\begin{itemize}
    \item \textbf{Imagen:} Es un archivo o plantilla de solo lectura que contiene las instrucciones, el código, el entorno de ejecución, las librerías y las variables necesarias para crear un contenedor.
    
    \item \textbf{Contenedor:} Es una instancia ejecutable de una imagen. Representa el proceso en ejecución aislado dentro del sistema anfitrión.
    
    \item \textbf{Volumen:} Es un mecanismo de almacenamiento persistente independiente del ciclo de vida del contenedor que permite guardar y compartir datos generados por las aplicaciones directamente en el sistema de archivos del anfitrión.
    
    \item \textbf{Puerto Publicado:} Es el mapeo explícito entre un puerto del sistema operativo anfitrión y un puerto interno del contenedor (\texttt{puerto\_anfitrión:puerto\_contenedor}), lo que permite dirigir el tráfico de red externo hacia el servicio que corre dentro del contenedor.
\end{itemize}

\subsection*{3. Importancia del Volumen y Persistencia de Datos}
El \textbf{volumen} es un elemento indispensable al desplegar un sistema gestor de bases de datos en un contenedor debido a la naturaleza efímera e inmutable de los contenedores. Por diseño, cuando un contenedor se detiene y se elimina, todas las modificaciones realizadas en su capa de almacenamiento de escritura se pierden de forma permanente. 

Si no se declara un volumen vinculado al directorio de datos del gestor (por ejemplo,\\ \texttt{/var/lib/postgresql/data} en PostgreSQL), todas las bases de datos, tablas y registros creados durante la ejecución del contenedor desaparecerán en el momento en que el contenedor sea removido o reiniciado. El volumen redirige la escritura hacia el almacenamiento físico de la computadora anfitriona, garantizando la persistencia de la información.

\section*{Fuentes Consultadas}
\begin{itemize}
    \item Docker Docs. (2026). \textit{Docker Overview \& Volumes}. https://docs.docker.com/
    \item PostgreSQL Global Development Group. (2026). \textit{PostgreSQL 16 Documentation}.\\ https://www.postgresql.org/docs/16/index.html
\end{itemize}

\end{document}